% BEGIN REV09_SUPERVIVENCIA_PUNTO_FIJO
\subsection{Survival by pruning prolongations}
\label{corpus9:xi:supervivencia}

Retrospective memory distinguishes histories recorded by the
APP--TRIT--TPK operations; prospective compatibility organizes their
admissible prolongations. Let \(S\) be the space of complete states
reachable through these operations. Each state retains sheets, residues,
quotients, orientation, carry, memory, and boundary data; whenever a
transition depends on depth, depth forms part of the state.
Let \(E\subseteq S\times S\) be the effective prolongation relation and
\(A\subseteq S\) the admissible set determined by the compatibility of
these records. Admissibility is thus decided through the operations and
boundary conditions that generate the history.

For \(X\subseteq S\), define
\begin{equation}
 \operatorname{Pre}_{E}(X)
 =\{s\in S:\exists t\in X,\ (s,t)\in E\},
 \qquad
 F_{\mathrm{surv}}(X)=A\cap\operatorname{Pre}_{E}(X).
 \label{corpus9:xi:operador-poda}
\end{equation}
The pruning sequence is given by
\begin{equation}
 V_0=A,\qquad
 V_{n+1}=F_{\mathrm{surv}}(V_n),\qquad
 V_\infty=\bigcap_{n\geq0}V_n.
 \label{corpus9:xi:sucesion-poda}
\end{equation}

\begin{proposition}[Survival and the greatest fixed point]
\label{corpus9:xi:mayor-punto-fijo}
The set \(V_n\) consists exactly of the states in \(A\) admitting an
\(n\)-step prolongation all of whose states belong to \(A\). The
sequence \((V_n)\) is decreasing. If each state has finitely many
admissible successors, then
\[
 V_\infty=F_{\mathrm{surv}}(V_\infty)
          =\nu F_{\mathrm{surv}},
\]
where \(\nu F_{\mathrm{surv}}\) denotes the greatest fixed point under
inclusion. Its elements are precisely the states admitting an infinite
compatible prolongation.
\end{proposition}

\begin{proof}
For \(n=0\), a prolongation of length zero consists of the state
itself in \(A\). If the characterization holds for \(n\), a state
belongs to \(V_{n+1}\) exactly when it belongs to \(A\) and has a
successor in \(V_n\); prefixing that edge produces an admissible
prolongation of \(n+1\) steps. This equivalence proves the assertion by
induction. Moreover, \(V_1\subseteq V_0\), and monotonicity of
\(F_{\mathrm{surv}}\) propagates \(V_{n+1}\subseteq V_n\).

Let \(s\in V_\infty\), and write
\(E(s)=\{t\in S:(s,t)\in E\}\). Since \(s\in V_{n+1}\) for every \(n\),
the sets \(E(s)\cap V_n\) are nonempty, finite, and decreasing.
Their intersection contains an element \(t\), which belongs to
\(V_\infty\) and satisfies \((s,t)\in E\). Hence
\(V_\infty\subseteq F_{\mathrm{surv}}(V_\infty)\).
Conversely, monotonicity gives
\(F_{\mathrm{surv}}(V_\infty)\subseteq V_{n+1}\) for every \(n\);
the image also lies in \(A=V_0\), and therefore belongs to
\(V_\infty\). This proves the fixed-point equality.

If \(W=F_{\mathrm{surv}}(W)\), then \(W\subseteq A=V_0\).
From \(W\subseteq V_n\) it follows that
\(W=F_{\mathrm{surv}}(W)\subseteq F_{\mathrm{surv}}(V_n)=V_{n+1}\).
Induction yields \(W\subseteq V_\infty\), proving maximality.
Finally, the prolongation tree of a state in \(V_\infty\) is finitely
branching and has a nonempty level at every depth. K\"onig's lemma
provides an infinite branch. Conversely, the prefixes of any infinite
prolongation place its initial state in every \(V_n\).
\end{proof}

\paragraph{Branching and stabilization.}
The finiteness hypothesis has an explicit control. Consider a root
with a distinct terminal branch of each positive integer length, all
branches being admissible and disjoint after the root. The root
admits prolongations of every finite length and belongs to all the
\(V_n\); each of its successors belongs to a branch of bounded length.
Consequently, the root has neither an infinite prolongation nor a
successor in \(V_\infty\). This tree has infinite branching at the
root. In contrast, if \(A\) is finite, each strict inclusion
\(V_{n+1}\subsetneq V_n\) removes at least one state. Pruning reaches
a fixed point after at most \(|A|\) strict layer removals.
This bound pertains to the finite domain \(A\).

\paragraph{Uniqueness of continuation.}
For \(s\in V_\infty\), the successors retaining survival form
\[
 E(s)\cap V_\infty.
\]
The next state is uniquely selected exactly when
\begin{equation}
 \#\bigl(E(s)\cap V_\infty\bigr)=1.
 \label{corpus9:xi:unicidad-sucesor}
\end{equation}
An infinite prolongation is unique when this condition holds at
each state of the path. Indeed, each surviving successor has an
infinite prolongation by the preceding proposition; two distinct
successors produce two distinct paths. When the successor is unique
at each step, the prefixes are determined successively.
Several viable successors express multiplicity of prolongations
within the same survival domain.

\begin{proposition}[Conjugation of the two directions of prolongation]
\label{corpus9:xi:conjugacion-supervivencia}
Let \(C:S\to S\) be an involution such that \(C(A)=A\) and
\begin{equation}
 (x,y)\in E\quad\Longleftrightarrow\quad(Cy,Cx)\in E.
 \label{corpus9:xi:involucion-prolongaciones}
\end{equation}
Let \(V_n^+\) denote pruning for \(E\), and \(V_n^-\) pruning for
\(E^{-1}\), both with initial set \(A\). Then
\[
 C(V_n^+)=V_n^-\quad(n\geq0),\qquad
 C(V_\infty^+)=V_\infty^-.
\]
\end{proposition}

\begin{proof}
If \(s_0,\ldots,s_n\) is an admissible path for \(E\), the sequence
\(Cs_0,\ldots,Cs_n\) is an admissible path for \(E^{-1}\):
for each \(j\), the hypothesis converts
\((s_j,s_{j+1})\in E\) into \((Cs_{j+1},Cs_j)\in E\).
All its states belong to \(A\). Applying the same operation again
recovers the initial path because \(C^2=I\); it therefore establishes
a bijection between finite prolongations in the two directions.
The characterization of \(V_n^\pm\) by paths gives
\(C(V_n^+)=V_n^-\). Since \(C\) is bijective, it commutes with the
intersection of these sets, yielding the equality for
\(V_\infty^\pm\).
\end{proof}

The involution transports the surviving set in one direction to
the surviving set in the conjugate direction. Infinite
prolongability in both directions corresponds to
\(V_\infty^+\cap V_\infty^-\); uniqueness of a path retains the
criterion in \eqref{corpus9:xi:unicidad-sucesor}. Retaining a
genealogy and restricting its continuation tree are two typed
operations on the same records. Assigning an entropy to these
readings additionally requires the corresponding measure on histories.
% END REV09_SUPERVIVENCIA_PUNTO_FIJO
